\documentclass[12pt, a4paper]{article}

\usepackage[T2A]{fontenc}
\usepackage[utf8]{inputenc}
\usepackage[russian]{babel}

\usepackage{geometry}
\geometry{a4paper, top=2cm, bottom=2cm, left=3cm, right=1.5cm}

\begin{document}

\begin{titlepage}
    \centering
    \vspace*{\stretch{0.5}}
    \textbf{Национальный исследовательский ядерный университет «МИФИ»} \\
    \vspace{0.5cm}
    Институт интеллектуальных кибернетических систем \\
    \vspace{0.5cm}
    Кафедра кибернетики (№22) \\
    \vspace{2cm}
    \textbf{\Large Курсовая работа} \\
    \vspace{0.5cm}
    \large по дисциплине "Современные технологии проектирования интеллектуальных систем" \\
    \vspace{2cm}
    \textbf{\huge Тема: «Интеллектуальная система поддержки принятия решений водителем при дорожно-транспортном происшествии»} \\
    \vspace{\stretch{1.5}}
    
    \begin{flushright}
        \begin{tabular}{l}
            \large
            \textbf{Выполнил:} \\
            студент группы М25-504 \\
            Сироджев А.Ш. \\
            \\
            \\
            \textbf{Научный руководитель:} \\
            д.т.н., проф. \\
            Рыбина Г.В. \\
        \end{tabular}
    \end{flushright}
    
    \vspace{\stretch{1}}
    \large Москва, 2025
\end{titlepage}

\section*{Проблемная область}

Разработка интеллектуальной системы-помощника для информационной поддержки водителя, попавшего в дорожно-транспортное происшествие (ДТП). Система предназначена для работы на мобильном устройстве и призвана снизить уровень стресса, минимизировать вероятность ошибок при оформлении документов и помочь водителю принять верные решения в критической ситуации, основываясь на анализе текущих обстоятельств.

\section*{Планируемые задачи}

В рамках курсовой работы предполагается решить следующий набор задач:

\begin{enumerate}
    \item \textbf{Сбор первичной информации о ДТП.}
    Система в режиме диалога (опросника) собирает ключевые данные о происшествии: наличие пострадавших, количество участников, расположение транспортных средств, наличие помех для движения.
    \textit{(Формализованная задача)}

    \item \textbf{Формирование динамического алгоритма действий для водителя.}
    Это ключевая задача системы. На основе анализа всей совокупности факторов (наличие пострадавших, согласие/несогласие второго участника, тип ДТП, погодные условия) система генерирует персонализированную пошаговую инструкцию. Алгоритм не является статичным и может меняться в зависимости от вводимых пользователем данных.
    \textit{\textbf{(Неформализованная задача - Планирование)}}

\end{enumerate}

\newpage
\section*{Дополнительный вариант №1: Интерпретация}

\textbf{Тема: «Интеллектуальная система для анализа и интерпретации отзывов пользователей мобильного приложения»}

\subsection*{Проблемная область}

Автоматизация обработки обратной связи от пользователей для выявления ключевых проблем и оценки удовлетворенности продуктом.

\subsection*{Планируемые задачи}

\begin{enumerate}
    \item \textbf{Агрегация отзывов.} Сбор текстовых отзывов из различных источников (Google Play, App Store, социальные сети) в единую базу данных.
    \textit{(Формализованная задача)}

    \item \textbf{Тематическая классификация.} Автоматическое определение основной темы отзыва (например, "проблемы с оплатой", "неудобный интерфейс", "запрос новой функции").
    \textit{(Формализованная задача)}

    \item \textbf{Интерпретация скрытого смысла и тональности отзыва.} Система анализирует не только прямые слова, но и контекст, чтобы определить суть проблемы, даже если она описана расплывчато (например, интерпретировать "после обновления всё сломалось" как проблему с совместимостью).
    \textit{\textbf{(Неформализованная задача - Интерпретация)}}

    \item \textbf{Формирование сводного отчета.} Создание аналитической сводки для разработчиков с указанием наиболее острых проблем, их частоты и критичности.
    \textit{(Формализованная задача)}
\end{enumerate}


\newpage
\section*{Дополнительный вариант №2: Диагностика}

\textbf{Тема: «Экспертная система для предварительной диагностики неисправностей бытовой техники»}

\subsection*{Проблемная область}

Создание интерактивного помощника, который помогает пользователю определить причину поломки бытового прибора (например, стиральной машины) и оценить возможность самостоятельного ремонта.

\subsection*{Планируемые задачи}
\begin{enumerate}
    \item \textbf{Сбор первичных "симптомов".} Пользователь выбирает устройство и описывает проблему с помощью предложенных вариантов ("не включается", "сильно шумит").
    \textit{(Формализованная задача)}

    \item \textbf{Выдача простых инструкций.} На основе первичного симптома система предлагает очевидные решения (например, "проверьте, включен ли прибор в розетку").
    \textit{(Формализованная задача)}

    \item \textbf{Проведение интерактивной диагностики.} Система запускает диалог, задавая уточняющие вопросы ("Шум появляется во время набора воды или при отжиме?"). Логика диалога имитирует рассуждения мастера по ремонту.
    \textit{\textbf{(Неформализованная задача - Диагностика)}}
    
    \item \textbf{Вынесение вердикта и рекомендаций.} По результатам диагностики система выдает заключение о наиболее вероятной причине поломки и предлагает план действий.
    \textit{(Формализованная задача)}
\end{enumerate}

\end{document}
